\section*{الفصل \ref{ch:b3:rate-theories} --- نظريات سرعة التفاعل}

\begin{solution}{exo:b3:rate-theories:1}
$\bar v_{\mathrm{rel}} = \sqrt{8 \times 1.381\times10^{-23} \times 298/(\pi \times 18.46 \times 1.661\times10^{-27})} = \qty{585}{m/s}$.
$\sigma\bar v_{\mathrm{rel}}N_A = 0.43\times10^{-18} \times 585 \times 6.022\times10^{23} = \qty{1.51e8}{m^3.mol^{-1}.s^{-1}} =
\qty{1.5e11}{L.mol^{-1}.s^{-1}}$. والمقيس: $2.07\times10^{-12} \times 6.022\times10^{23}\times10^{-3} =
\qty{1.25e9}{L.mol^{-1}.s^{-1}}$، ومنه $P = 0.008$: فأقل من تصادم واحد من كل مئة له
الاتجاه الصحيح (يجب أن تلاقي ذرة O في NO ذرة O طرفية في \ce{O3}).
\end{solution}

\begin{solution}{exo:b3:rate-theories:2}
$\Delta^{\ddagger}H^\circ = 75.0 - 2.48 = \qty{72.5}{kJ/mol}$ في المحلول؛ و$75.0 - 4.96 = \qty{70.0}{kJ/mol}$
لتفاعل غازي ثنائي الجزيئية.
\end{solution}

\begin{solution}{exo:b3:rate-theories:3}
ديلز--ألدر: سالبة بشدة (جزيئان يكوّنان معقدًا حلقيًا واحدًا).
فقد \ce{N2}: موجبة (رابطة ترتخي، وتُكتسب حرية). تولّد أيونات في
مذيب قطبي: سالبة، فالشحنات الناشئة تنظّم المذيب من حولها.
\end{solution}

\begin{solution}{exo:b3:rate-theories:4}
الماء: $8 \times 8.314 \times 298.15/(3 \times \num{0.890e-3}) = \qty{7.4e9}{L.mol^{-1}.s^{-1}}$؛ والهكسان
\qty{2.2e10}{L.mol^{-1}.s^{-1}}. الهكسان، الأقل لزوجة بثلاث مرات: فالجزيئات تنتشر فيه أسرع بثلاث
مرات.
\end{solution}

\begin{solution}{exo:b3:rate-theories:5}
$\ln(k_2T_1/k_1T_2) = \ln(10 \times 298.15/318.15) = 2.238$؛ $\Delta^{\ddagger}H^\circ = 8.3145 \times 2.238/(1/298.15 -
1/318.15) = \qty{88.2}{kJ/mol}$. $\Delta^{\ddagger}S^\circ = R[\ln(kh/k_BT) + \Delta^{\ddagger}H^\circ/RT] = 8.3145 \times
(-36.365 + 35.596) = \qty{-6.4}{J.K^{-1}.mol^{-1}}$ (عند \qty{298.15}{K}). $E_a = R\ln10/(1/298.15 - 1/318.15) =
\qty{90.8}{kJ/mol} \approx \Delta^{\ddagger}H^\circ + RT$ عند درجة الحرارة المتوسطة.
\end{solution}

\begin{solution}{exo:b3:rate-theories:6}
$\mu_{\mathrm{OH}} = 16/17$، $\mu_{\mathrm{OD}} = 32/18$، ونسبة العددين الموجيين $\sqrt{1.889} = 1.374$:
$\tilde\nu_{\mathrm{OD}} = \qty{2661}{cm^{-1}}$. $k_{\mathrm H}/k_{\mathrm D} = \exp(1.4388 \times 996/(2 \times 298)) = \eu^{2.40} = 11$.
\end{solution}

\begin{solution}{exo:b3:rate-theories:7}
القيمة العظمى $\exp(1.4388 \times 808/(2 \times 298)) = \eu^{1.95} = 7.0$. وقيمة 40 لا يمكن أن تأتي من
طاقات النقطة الصفرية: فالهيدروجين يعبر الحاجز نفقيًا، وهو ما يفعله الدوتيريوم،
الأثقل بمرتين، أقل بكثير.
\end{solution}

\begin{solution}{exo:b3:rate-theories:8}
$\ce{S2O8^{2-}}$ و\ce{I-}: $z_{\mathrm A}z_{\mathrm B} = +2$، $\log(k/k_0) = 2 \times 0.51 \times 2 \times 0.10 = 0.20$: يُضرب $k$
في 1.6. الإستر و\ce{OH-}: أحد الشريكين متعادل، فلا أثر.
$\ce{[Co(NH3)5Br]^{2+}}$ و\ce{OH-}: $z_{\mathrm A}z_{\mathrm B} = -2$، يُضرب $k$ في $10^{-0.20} = 0.63$.
\end{solution}

\begin{solution}{exo:b3:rate-theories:9}
مبكر، في وادي الدخول (حالة انتقالية شبيهة بالمتفاعلات، بحسب مسلّمة
هاموند). والتفاعل العكسي، الماص للحرارة، له حاجز متأخر على
طريقه الخاص: فالطاقة الاهتزازية في الرابطة H--F تساعده أكثر.
\end{solution}

\begin{solution}{exo:b3:rate-theories:10}
$\Delta^{\ddagger}H^\circ = -Rb = \qty{61.9}{kJ/mol}$، بخطأ معياري $R \times 31 = \qty{0.26}{kJ/mol}$ 
ومجال 95\,\% قدره $\pm\qty{0.8}{kJ/mol}$.
\[
  \Delta^{\ddagger}S^\circ = R(16.528 - 23.760) = \qty{-60.1}{J.K^{-1}.mol^{-1}},
\]
بخطأ معياري $R \times 0.103 = \qty{0.86}{J.K^{-1}.mol^{-1}}$ ومجال 95\,\% قدره $\pm\qty{2.7}{J.K^{-1}.mol^{-1}}$.
\end{solution}

\begin{solution}{exo:b3:rate-theories:11}
انطلاقًا من $\ln k = \ln(k_BT/h) + \Delta^{\ddagger}S^\circ/R - \Delta^{\ddagger}H^\circ/RT$،
\[
  \frac{\dd\ln k}{\dd T} = \frac1T + \frac{\Delta^{\ddagger}H^\circ}{RT^2}, \qquad
  E_a = RT^2\frac{\dd\ln k}{\dd T} = \Delta^{\ddagger}H^\circ + RT .
\] ثم $A = k\eu^{E_a/RT} = (k_BT/h)\eu^{\Delta^{\ddagger}S^\circ/R}\eu^{-\Delta^{\ddagger}H^\circ/RT}\eu^{(\Delta^{\ddagger}H^\circ + RT)/RT} =
\eu\,(k_BT/h)\eu^{\Delta^{\ddagger}S^\circ/R}$. ومع $\Delta^{\ddagger}S^\circ = 0$: $A = 2.718 \times \num{6.21e12} = \qty{1.7e13}{s^{-1}}$، وهي
رتبة العوامل المرصودة؛ والقيم التي تفوقها تعني $\Delta^{\ddagger}S^\circ$ موجبة، أي حالة
انتقالية رخوة.
\end{solution}

\begin{solution}{exo:b3:rate-theories:12}
مع دوال تجزئة جزيئية لوحدة الحجم: $q^{\mathrm T}/V = (2\pi mk_BT)^{3/2}/h^3$ من أجل
A وB والمعقد (الكتلة $m_{\mathrm A} + m_{\mathrm B}$)، ومن أجل المعقد $q^{\mathrm R} = 8\pi^2Ik_BT/h^2$
($\sigma = 1$). والنسبة الانسحابية $h^3/(2\pi\mu k_BT)^{3/2}$، ومنه
\[
  k = \frac{k_BT}{h}\cdot\frac{h^3}{(2\pi\mu k_BT)^{3/2}}\cdot\frac{8\pi^2\mu d^2k_BT}{h^2}\,\eu^{-E_0/RT}
    = \frac{8\pi^2d^2(k_BT)^{1/2}}{(2\pi)^{3/2}\mu^{1/2}}\eu^{-E_0/RT}
    = \pi d^2\sqrt{\frac{8k_BT}{\pi\mu}}\,\eu^{-E_0/RT},
\]
لكل زوج من الجزيئات؛ والضرب في $N_A$ يعطي \cref{thm:b3:rate-theories:collision}
مع $\sigma = \pi d^2$.
\end{solution}

\begin{solution}{pb:b3:rate-theories:1}
\textbf{1.} $\mu_{\mathrm{CH}} = 12 \times 1.00783/13.00783 = \qty{0.9297}{u}$؛ $\mu_{\mathrm{CD}} = 12 \times 2.01410/14.01410 =
\qty{1.7246}{u}$.
\textbf{2.} $2917/\sqrt{1.7246/0.9297} = 2917/1.3620 = \qty{2142}{cm^{-1}}$، أي أعلى بنسبة 1.5\,\% من
القيمة المقيسة 2109 (اللاتوافق واقتران الروابط الأربع في \ce{CD4}).
\textbf{3.} $\frac12 \times 2917 = \qty{1458.5}{cm^{-1}}$ و$\frac12 \times 2109 = \qty{1054.5}{cm^{-1}}$؛ والفرق
\qty{404}{cm^{-1}} $= \qty{4.83}{kJ/mol}$.
\textbf{4.} يصير إحداثي التفاعل: فتختفي طاقة نقطته الصفرية.
\textbf{5.} الحاجز أعلى من أجل D بمقدار \qty{4.83}{kJ/mol}.
\textbf{6.} $\exp(1.43878 \times 404/310.15) = \eu^{1.874} = 6.5$.
\textbf{7.} $\exp(1.43878 \times 404/298.15) = 7.0$.
\textbf{8.} في حالة انتقالية حقيقية يبقى جزء من \omterm{def:b3:quantum-model-systems:oscillator}{طاقة النقطة الصفرية} للتمدد
في أنماط أخرى، وهذا يُضعف المفعول؛ ولا تستطيع طاقات النقطة الصفرية وحدها
أن تتجاوز هذه القيمة (أما \omterm{def:b3:quantum-model-systems:tunnelling}{العبور النفقي} فيستطيع).
\textbf{9.} قليلًا: فهي مفاعيل ثانوية، من رتبة 10\,\% لكل دوتيريوم على
الأكثر.
\textbf{10.} أن الرابطة C--H تنكسر في الخطوة المحددة للسرعة، مع
تحول التمدد كليًا إلى إحداثي التفاعل.
\textbf{11.} طرح من الرتبة الأولى: $t_{1/2} = \ln2/k_{\mathrm{total}}$، متناسب عكسيًا معه.
\textbf{12.} $k_{\mathrm{total}}(\mathrm H) = 0.75\,k_{\mathrm{total}}(\mathrm H) + 0.25\,k_{\mathrm{total}}(\mathrm H)$، نزع الميثيل والطرق الأخرى.
\textbf{13.} مقسوم على 6.5.
\textbf{14.} $0.25 + 0.75/6.52 = 0.365$.
\textbf{15.} $5.0/0.365 = \qty{13.7}{h}$.
\textbf{16.} 2.7.
\textbf{17.} 6.5، أي المفعول النظائري كاملًا.
\textbf{18.} جرعات أصغر أو أقل تواترًا، ومستويات أثبت في الدم.
\textbf{19.} $\ln(k/T)$ بدلالة $1/T$.
\textbf{20.} $\Delta^{\ddagger}H^\circ = 61.9 \pm \qty{0.3}{kJ/mol}$، $\Delta^{\ddagger}S^\circ = -60.1 \pm \qty{0.9}{J.K^{-1}.mol^{-1}}$.
\textbf{21.} $\Delta^{\ddagger}H^\circ = 66.8 \pm \qty{0.2}{kJ/mol}$، $\Delta^{\ddagger}S^\circ = -60.3 \pm \qty{0.7}{J.K^{-1}.mol^{-1}}$.
\textbf{22.} $\Delta\Delta^{\ddagger}H^\circ = 4.88 \pm \qty{0.33}{kJ/mol}$ ($\sqrt{0.26^2 + 0.21^2}$).
\textbf{23.} نعم: 4.83 تقع ضمن خطأ معياري واحد بفارق كبير. فالمفعول النظائري
تفسّره طاقات النقطة الصفرية كله، دون أي أثر \omterm{def:b3:quantum-model-systems:tunnelling}{للعبور النفقي}.
\textbf{24.} إنهما متساويتان في حدود خطأيهما: فللحالة الانتقالية البنية نفسها
\omterm{def:b3:rovibrational-spectroscopy:rotational-constant}{للنظيرين الجزيئيين} كليهما، كما افتُرض في الجزء الأول.
\textbf{25.} \textbf{$t_{1/2}(\mathrm D)/t_{1/2}(\mathrm H) \approx 2.7$ عند \qty{37}{\celsius}}: \qty{13.7}{h} مقابل \qty{5.0}{h}.
\end{solution}
